\subsection{Equilibrium Points and Operating Region Analysis}

When we linearize a nonlinear system around an operating point, we implicitly assume that the system will remain close to that point during operation. This assumption is not always valid, and understanding when it holds requires us to examine the underlying structure of the nonlinear dynamics more carefully. The key concepts that govern this analysis are equilibrium points, regions of attraction, and the global behavior captured by phase portraits. These tools allow us to determine not just where to linearize, but whether our linear approximation will be meaningful for the control tasks we have in mind.

An equilibrium point of a dynamical system is a state where the system can remain at rest indefinitely. Mathematically, a state vector $\mathbf{x}_e$ is an equilibrium point if when the system reaches this state, the derivatives of all state variables become identically zero. For the general nonlinear system $\dot{\mathbf{x}} = \mathbf{f}(\mathbf{x}, \mathbf{u})$, an equilibrium point corresponding to a constant input $\mathbf{u}_e$ satisfies the algebraic equation $\mathbf{0} = \mathbf{f}(\mathbf{x}_e, \mathbf{u}_e$. This condition means that at equilibrium, whatever driving forces exist in the system are perfectly balanced by whatever resisting forces exist, resulting in no change in the system state over time. Physical systems have equilibrium points whenever there exists a configuration where all accelerations vanish and all flows balance.

Finding equilibrium points typically requires solving a system of nonlinear algebraic equations. For simple systems with one or two states, we can often solve these equations analytically by inspection. For more complex systems, we may need to use numerical methods such as Newton-Raphson iteration. The number of equilibrium points a system possesses depends critically on its nonlinear structure. Some systems have a single equilibrium point that attracts all possible initial conditions, while others have multiple equilibrium points separated by boundaries that the system cannot cross without external intervention. Understanding which equilibrium points are reachable and stable is fundamental to selecting appropriate operating points for control.

Consider the simple pendulum system to illustrate equilibrium point analysis. The dynamics of a pendulum of length $L$ with mass $m$ and gravitational acceleration $g$ are given by $ml^2\ddot{\theta} + b\dot{\theta} + mgL\sin\theta = u$, where $\theta$ is the angular displacement from vertical and $u$ is an applied torque. Defining state variables $x_1 = \theta$ and $x_2 = \dot{\theta}$, we obtain the state-space representation $\dot{x}_1 = x_2$ and $\dot{x}_2 = -\frac{g}{L}\sin x_1 - \frac{b}{ml^2}x_2 + \frac{u}{ml^2}$. To find equilibrium points, we set the time derivatives to zero, which gives us $0 = x_2$ and $0 = -\frac{g}{L}\sin x_1 + \frac{u}{ml^2}$. From the first equation, we immediately see that at equilibrium the angular velocity must be zero. The second equation then requires that $-\frac{g}{L}\sin x_1 + \frac{u}{ml^2} = 0$, which we can solve for the equilibrium angles as $\sin x_1 = \frac{Lu}{mgL^2} = \frac{u}{mgL}$. Since the sine function is bounded between $-1$ and $1$, equilibrium points exist only when the applied torque satisfies $|u| \leq mgL$. When $u = 0$, we have $\sin x_1 = 0$, giving equilibrium points at $\theta = 0$ (hanging down) and $\theta = \pi$ (pointing straight up). When $u$ is positive, the equilibrium shifts downward in the direction of the applied torque. This simple example shows how equilibrium points depend on both the system state and the constant input we apply.

Once we identify equilibrium points, we must assess their stability before selecting an operating point for control. A stable equilibrium point is one where, if the system is perturbed slightly away from equilibrium, it will naturally return to equilibrium without external intervention. An unstable equilibrium point is one where any perturbation, no matter how small, causes the system to move away from equilibrium. For control purposes, we almost always want to operate near a stable equilibrium point because this gives us a natural restoring tendency that our controller can exploit. Operating near an unstable equilibrium requires continuous active control to maintain the desired state, which is more demanding and offers less margin for error.

Linearization around an equilibrium point reveals stability properties through the eigenvalues of the Jacobian matrix evaluated at that point. For the pendulum example with $u = 0$, we linearize about the hanging-down equilibrium $\theta_e = 0$. The Jacobian matrix at this equilibrium is $A = \begin{bmatrix} 0 & 1 \\ -\frac{g}{L} & -\frac{b}{ml^2} \end{bmatrix}$. The characteristic equation is $\det(sI - A) = s^2 + \frac{b}{ml^2}s + \frac{g}{L} = 0$. For any positive damping coefficient $b > 0$, both eigenvalues have negative real parts, confirming that the hanging-down equilibrium is stable. In contrast, linearizing about the upright equilibrium $\theta_e = \pi$ gives $A = \begin{bmatrix} 0 & 1 \\ \frac{g}{L} & -\frac{b}{ml^2} \end{bmatrix}$, with characteristic equation $s^2 + \frac{b}{ml^2}s - \frac{g}{L} = 0$. This equation always has at least one eigenvalue with positive real part, indicating that the upright equilibrium is unstable regardless of how much damping we add. This is a fundamental property of the inverted pendulum that no linear controller can change, though we can stabilize it using nonlinear or switching control techniques.

Phase portraits provide a powerful visual tool for understanding the global behavior of nonlinear systems. A phase portrait is a plot in the state space showing trajectories that the system would follow from various initial conditions, along with arrows indicating the direction of motion. For second-order systems with states $x_1$ and $x_2$, we plot $x_2$ versus $x_1$ and draw curves showing how the system evolves from different starting points. Equilibrium points appear as special points in this diagram, and their stability determines the local geometry of nearby trajectories. Stable equilibria appear as focal points where trajectories spiral inward or as nodal points where trajectories approach directly. Unstable equilibria appear as sources (trajectories moving away) or saddle points (trajectories approaching along some directions but departing along others).

The pendulum phase portrait beautifully illustrates the different behaviors possible in a nonlinear system. The hanging equilibrium at the origin appears as a spiral sink, meaning that all trajectories sufficiently close to the origin spiral inward and eventually approach the equilibrium. The upright equilibrium at $\theta = \pi$ appears as a saddle point, with trajectories approaching along the vertical direction but departing along horizontal directions. Separating these two regions of qualitatively different behavior is a curve called the separatrix, which passes through the saddle point. Trajectories that lie inside the region bounded by the separatrix eventually approach the stable hanging equilibrium, while trajectories outside this region spiral outward and represent pendulum swings that complete full rotations. The separatrix itself represents the boundary case of a pendulum that approaches but never quite reaches the upright position.

The region from which a system approaches a particular stable equilibrium is called the region of attraction or basin of attraction. For the pendulum with zero applied torque, the region of attraction for the hanging equilibrium is the set of all initial conditions where the total mechanical energy is less than the potential energy at the upright position. Mathematically, this corresponds to $\frac{1}{2}ml^2x_2^2 + mgL(1 - \cos x_1) < 2mgL$. States on the separatrix have exactly enough energy to reach the upright position with zero velocity, and any slight perturbation can push them into the region that leads to full rotation. This concept of regions of attraction is absolutely essential for control design because it tells us whether our linearized model will be valid during actual operation. If our controller needs to move the system through regions where the linearized dynamics are poor approximations, we may encounter unexpected behavior that the linear design did not anticipate.

Consider what happens when we linearize the pendulum only about the hanging equilibrium and then attempt to swing the pendulum up to the upright position using this linear controller. The linearized dynamics near the hanging position suggest that we can achieve good tracking performance by applying appropriate torques. However, as the pendulum swings upward, it enters regions of the state space where the linear approximation becomes increasingly inaccurate. Near the upright position, the restoring torque reverses direction compared to what the linear model predicts. A controller designed purely from the linear model will therefore be surprised by the system's actual behavior and may fail to stabilize at the upright position, or worse, may cause the pendulum to spin out of control. This example demonstrates that understanding the validity range of our linearized model is not an academic exercise but a practical necessity for successful control design.

The selection of an operating point for linearization and control design should be guided by several practical considerations. First, we must choose an equilibrium point that lies within a region where the system will actually spend most of its operating time. If our control objective requires the system to move between different regions of the state space, we may need multiple linearized models at different operating points, or we may need to use a controller designed directly from the nonlinear model. Second, we must assess whether the dynamics near the chosen operating point are representative of the dynamics that matter for our control objectives. Some equilibrium points are highly sensitive to perturbations, meaning that even small disturbances will move the system into regions where the linear approximation fails. Third, we must consider whether the linearized model captures the coupling between states that our controller needs to manage. Linearization can sometimes mask important nonlinear interactions that become significant during transient responses.

The concept of operating region analysis extends beyond single equilibrium points to consider collections of states where a particular linearized model provides adequate accuracy. We typically define an operating region as an ellipsoidal or rectangular region in the state space centered on the equilibrium point, sized such that the error between the nonlinear system and its linear approximation remains below some acceptable threshold. This threshold might be specified in terms of state errors, output errors, or control effort differences. The size and shape of the valid operating region depend on the specific nonlinearities present in the system and on how aggressive our accuracy requirements are. For systems with mild nonlinearities, the operating region might be quite large, allowing us to use a single linear controller over a wide range of operating conditions. For systems with strong nonlinearities, the operating region might be quite small, requiring gain scheduling or other techniques to maintain performance across the full operating envelope.

In practice, operating region analysis often proceeds through a combination of analytical bounds and numerical simulation. We can derive analytical bounds on linearization error using Taylor's theorem with remainder, which gives us an explicit expression for how the approximation error grows with distance from the operating point. For a scalar system $\dot{x} = f(x)$, the error between the actual nonlinear dynamics and the linearized dynamics $\dot{x} = a x$ where $a = f'(x_e)$ satisfies $|\dot{x} - a x| \leq \frac{1}{2}M(x - x_e)^2$ where $M$ is the maximum absolute value of the second derivative $f''(x)$ in the region of interest. This bound tells us that the linearization error grows quadratically with distance from the operating point, which explains why linear models are most accurate near equilibrium and become increasingly inaccurate far from equilibrium. For multi-dimensional systems, similar bounds involve the Hessian matrices of the vector field components.

Numerical simulation provides a more precise but less general approach to characterizing operating regions. We can simulate the closed-loop system with our linear controller from many different initial conditions and observe whether the performance remains acceptable. By systematically exploring the state space, we can map out the region where our controller achieves the desired performance specifications. This empirical characterization is particularly valuable when analytical bounds are too conservative or when the system dynamics are too complex to allow simple analytical expressions. The resulting region can be stored as a lookup table or approximated by a polynomial, and then used by a higher-level supervisory controller to determine whether gain scheduling or model adaptation is needed.

The interaction between operating point selection and control design decisions is subtle and important. A control system designed to regulate a particular equilibrium point will typically perform best when operating near that point and may exhibit degraded performance or even instability when pushed too far away. This observation motivates a hierarchical control structure where a high-level supervisor selects among multiple low-level controllers, each designed for a different operating region. Alternatively, we can design a single controller that maintains validity across a larger operating region by explicitly accounting for the nonlinear dynamics in the controller design. This might involve feedback linearization, sliding mode control, or other nonlinear control techniques that achieve global validity at the cost of increased complexity in analysis and implementation.

In summary, the analysis of equilibrium points and operating regions provides essential context for any control design based on linearization. We must identify which equilibrium points are stable and reachable, understand the regions from which the system will approach these equilibria, assess how far we can push the system while maintaining adequate linear model accuracy, and select operating points that support our control objectives. This analysis transforms linear control design from a purely mathematical exercise into an engineering discipline that accounts for the practical realities of nonlinear system behavior. When we complete this analysis, we understand not just where to linearize, but whether our linear controller will continue to perform well under the conditions it will actually encounter in deployment.
